\documentclass[conference]{IEEEtran}
\usepackage{amsmath,amssymb}
\usepackage{algorithm}
\usepackage{algpseudocode}
\usepackage{booktabs}
\usepackage{graphicx}
\usepackage{textcomp}
\usepackage{xcolor}
\usepackage{url}

% Placeholder for values that come from Colab runs (see PROGRESS.md). Replace every occurrence
% before submission; never type a number that is not in results/*.json.
\newcommand{\pending}[1]{\textcolor{red}{TODO(user): #1}}

% Insert a figure when it exists, otherwise draw a labelled placeholder. This keeps the skeleton
% compilable on Overleaf before the long jobs have produced the PDFs.
\newcommand{\pendingfigure}[2]{%
  \IfFileExists{#1}%
    {\includegraphics[width=\linewidth]{#1}}%
    {\fbox{\parbox[c][0.22\textheight][c]{0.92\linewidth}{\centering\small #2}}}%
}

\begin{document}

\title{Criterion-Level Clinical Trial Screening with a Fine-Tuned Qwen LoRA:\\
KV-Cached Decoding, Tiled Attention and Speculative Decoding}

\author{\IEEEauthorblockN{Rayan Khan}
\IEEEauthorblockA{GenAI Assignment 2\\Email: rayankhannexus@gmail.com}}

\maketitle

\begin{abstract}
We study criterion-level patient-to-trial screening with a small instruction-tuned language model.
Given a patient note and one eligibility criterion, the model returns a label from the criterion's
legal label set together with the patient-sentence identifiers that support the decision; the
explanations shown to users are generated deterministically from the label and the cited sentences,
never by the model. We fine-tune Qwen2.5-1.5B-Instruct with a low-rank adapter on the TrialGPT
criterion annotations~\cite{jin2024trialgpt} and evaluate classification and evidence quality on a
patient-level split: validation accuracy rises from $1.5\%$ (base) to $70.4\%$ (fine-tuned), test
accuracy from $1.9\%$ to $71.4\%$, while macro-F1 stays near $0.37$ because the adapter never emits
the three rarest labels. The KV-cached decoder reproduces the Hugging Face baseline byte-for-byte on
all 206 validation records. On top of the fine-tuned model we implement three inference components from
scratch: an explicit KV-cache decoding loop with a loop-free reference, a tiled attention kernel
with an online softmax~\cite{milakov2018online,dao2022flashattention}, and greedy speculative
decoding with a 0.5B draft model~\cite{leviathan2023speculative,chen2023speculative}. All three are
tested for exact equivalence against their references and benchmarked on a T4 GPU. The KV-cached
decoder is token-identical to the framework baseline and $1.0$--$5.1\times$ faster; the tiled kernel
matches the reference within fp16 tolerance but is slower than fused SDPA while using six times less
memory at 2{,}048 tokens; greedy speculative decoding is lossless yet $1.5\times$ slower on this GPU
despite $63$--$82\%$ acceptance. A Gradio application exposes screening, report export and a live
decoder comparison.
\end{abstract}

\begin{IEEEkeywords}
clinical trial screening, parameter-efficient fine-tuning, LoRA, KV cache, tiled attention,
speculative decoding, large language models.
\end{IEEEkeywords}

\section{Introduction}
Recruiting participants for clinical trials is slow and expensive, and a large share of the effort
is spent reading free-text patient notes against long lists of eligibility criteria. Large language
models (LLMs) can draft criterion-level judgements, but three practical questions remain: how well a
small locally runnable model can do after parameter-efficient fine-tuning, whether its decisions can
be audited through cited evidence, and whether the inference tricks used in production systems are
worth their complexity at this scale.

This paper answers those questions for a single-criterion formulation. A note and one criterion are
rendered into a chat prompt; the model must reply with one JSON object
\texttt{\{"label": ..., "evidence\_ids": [...]\}}. The label set depends on whether the criterion is
an inclusion or an exclusion criterion, and evidence identifiers refer to the numbered sentences of
the note exactly as they appear in the dataset. GPT-4 annotations that accompany the dataset are
discarded at load time, so no distillation signal from the proprietary model is used.

Our contributions are:
\begin{itemize}
  \item a Qwen2.5-1.5B LoRA model for criterion-level screening, trained with completion-only loss
        on 603 examples and evaluated on a held-out patient-level split;
  \item a from-scratch KV-cached decoder, verified token-identical to an uncached reference, and a
        wall-clock study of cache, attention and speculative decoding on a T4 GPU;
  \item a tiled attention implementation with a running softmax maximum and normaliser, checked
        against a full-matrix reference and PyTorch SDPA;
  \item a lossless greedy speculative decoder with a 0.5B draft model that reports acceptance rate
        and the resulting speedup \emph{or slowdown};
  \item an application with a persistent educational disclaimer, deterministic explanations,
        downloadable reports and a live decoder-comparison tab.
\end{itemize}

\section{Related Work}
TrialGPT~\cite{jin2024trialgpt} introduced an end-to-end zero-shot pipeline for patient-to-trial
matching and released the criterion annotations used here; its matching module reports 87.3\%
accuracy on 1{,}015 patient--criterion pairs with human evaluation. We keep the criterion-level
formulation but fine-tune a much smaller open model and constrain the output to a compact JSON
schema so that evidence can be verified automatically.

LoRA~\cite{hu2022lora} freezes the base weights and learns low-rank updates, which makes fine-tuning
a 1.5B model feasible on a single 16\,GB GPU. The base model is Qwen2.5-1.5B-Instruct from the
Qwen2.5 family~\cite{yang2024qwen25}, whose 0.5B sibling serves as the speculative draft.

Our inference work builds on the transformer architecture~\cite{vaswani2017attention}, on
FlashAttention's tiling idea~\cite{dao2022flashattention} and on the online normaliser used inside
it~\cite{milakov2018online}. Speculative decoding was proposed independently by
Leviathan et al.~\cite{leviathan2023speculative} and Chen et al.~\cite{chen2023speculative}; both
preserve the target distribution, and the greedy variant we implement preserves the target's exact
token sequence.

\section{Materials and Methods}

\subsection{Dataset and splits}
We use \texttt{ncbi/TrialGPT-Criterion-Annotations}, which contains 1{,}015 annotated
patient--criterion pairs from 53 patients and 103 trials. Each row carries the patient note, the
criterion type and text, the expert label and the expert evidence sentence identifiers. Notes are
pre-numbered from zero; we verified that numbering is consecutive on all 53 patients. The four
legal labels per criterion type and their counts are summarised in Table~\ref{tab:data}.

Splitting is performed at the \emph{patient} level with seed 42, giving 37/8/8 patients and
603/206/206 rows for training, validation and test. Patient-level splitting prevents the same note
from appearing in two splits. The test split is evaluated exactly once, after all model selection,
through a guarded \texttt{-{}-final} flag.

\begin{table}[t]
\centering
\caption{Dataset summary and expert label distribution. Counts are computed from the loader; the
split is frozen in \texttt{splits/split\_manifest.json}.}
\label{tab:data}
\begin{tabular}{lr}
\toprule
Property & Value \\
\midrule
Annotated pairs & 1{,}015 \\
Patients / trials & 53 / 103 \\
Train / validation / test rows & 603 / 206 / 206 \\
\midrule
Inclusion: included & 150 \\
Inclusion: not included & 23 \\
Inclusion: not enough information & 174 \\
Inclusion: not applicable & 9 \\
Exclusion: excluded & 10 \\
Exclusion: not excluded & 484 \\
Exclusion: not enough information & 119 \\
Exclusion: not applicable & 46 \\
\bottomrule
\end{tabular}
\end{table}

\subsection{Task, prompt and target}
For one criterion we render the numbered note sentences, the criterion type and text, and the legal
label set into a user message; a fixed system message states that the answer must be one JSON
object. The supervised target is compact JSON with sorted evidence identifiers,
e.g.\ \texttt{\{"label": "included", "evidence\_ids": [2, 5]\}}. Tokens belonging to the prompt are
masked out of the loss, so the model is trained only on its own completion.

Evidence identifiers outside the note are dropped by the parser, and an answer whose label is not in
the legal set for that criterion type is recorded as \texttt{invalid} and counted as wrong in every
metric. The parser first attempts a strict JSON decode, then recovers the largest brace-delimited
object, then falls back to regex extraction; repaired answers are flagged so the error analysis can
separate formatting failures from reasoning failures.

\subsection{Base model and fine-tuning}
The base model is Qwen2.5-1.5B-Instruct. We attach a LoRA adapter to all attention and MLP
projections and train for three epochs with the hyper-parameters in Table~\ref{tab:hparams}. The
only preprocessing beyond chat templating is truncation at 2{,}048 tokens; a separate analysis
showed that no training example exceeds 595 prompt-plus-target tokens, so no supervision is lost.
Training logs and curves are written to \texttt{results/train\_log.json} and
\texttt{report/figures/train\_curves.pdf}.

\begin{table}[t]
\centering
\caption{LoRA fine-tuning hyper-parameters (values fixed in \texttt{constants.py}).}
\label{tab:hparams}
\begin{tabular}{ll}
\toprule
Parameter & Value \\
\midrule
LoRA rank / alpha / dropout & 16 / 32 / 0.05 \\
Target modules & all attention and MLP projections \\
Learning rate / schedule & $2\times10^{-4}$ / cosine \\
Epochs / batch / accumulation & 3 / 2 / 8 \\
Weight decay / grad clip & 0.01 / 1.0 \\
Precision & fp16 on CUDA \\
\bottomrule
\end{tabular}
\end{table}

\subsection{Screening pipeline and aggregation}
For a new note and a list of criteria, the pipeline renders one prompt per criterion and calls the
same greedy decoder used in evaluation. Each prediction becomes a row with the criterion, its type,
the predicted label, the evidence identifiers, the quoted evidence sentences and a deterministic
explanation template. A separate aggregator produces a \emph{preliminary} patient-level outcome:
\texttt{not applicable} criteria are ignored; any inclusion \texttt{not included} or exclusion
\texttt{excluded} gives \emph{Likely ineligible}; any \texttt{not enough information} or
unparseable label gives \emph{Needs more information}; and only when every inclusion is
\texttt{included} and every exclusion is \texttt{not excluded} is the result \emph{Likely eligible}.
This rule is stated in the exported report; it is a heuristic, not a clinical decision.

\subsection{KV-cached greedy decoding}
The reference decoder re-runs the full sequence at every step with the cache disabled. The cached
decoder performs one prefill and then feeds one token per step while growing a
\texttt{DynamicCache}. Because rotary positions are absolute, the incremental step passes the cache
length as its position; the attention mask covers every cached position plus the new one. Greedy
decoding makes the two paths token-identical in exact arithmetic; in fp16 the reduced shapes can
change the last bits, so the benchmark reports any real-model mismatches instead of assuming
equality. Algorithm~\ref{alg:cached} states the loop.

\begin{algorithm}[t]
\caption{KV-cached greedy decoding}
\label{alg:cached}
\begin{algorithmic}[1]
\Require prompt tokens $x_{1..L}$, budget $N$
\State $c \gets \emptyset$; $y \gets []$
\State $(c, \ell) \gets \textsc{Forward}(x_{1..L}, c, \text{positions}=1..L)$
\For{$i \gets 1$ to $N$}
  \State $t \gets \arg\max \ell$; append $t$ to $y$
  \If{$t$ is EOS} \State \textbf{break} \EndIf
  \State $n \gets \text{length}(c)$
  \State $(c, \ell) \gets \textsc{Forward}([t], c, \text{positions}=n)$
\EndFor
\State \Return $y$
\end{algorithmic}
\end{algorithm}

\subsection{Tiled attention with an online softmax}
The reference attention materialises the full $T\times T$ score matrix and normalises it with
\texttt{softmax}. The tiled implementation processes query and key blocks and maintains, per query
row, a running maximum $m$, a running normaliser $\ell$ and a running accumulator $a$. When a new
key block raises the maximum, the accumulator and normaliser are rescaled by
$\exp(m_{\text{old}}-m_{\text{new}})$ before the block's contribution is added, so the final result
equals the full softmax without ever storing it (Algorithm~\ref{alg:tiled}). Causal masking is
applied per block, and fully masked rows keep a zero denominator so they yield zeros rather than
NaNs.

\begin{algorithm}[t]
\caption{Tiled attention with online softmax (one query block)}
\label{alg:tiled}
\begin{algorithmic}[1]
\Require $Q \in \mathbb{R}^{B\times H\times T\times D}$, $K,V \in \mathbb{R}^{B\times H\times S\times D}$
\State $m \gets -\infty$, $\ell \gets 0$, $a \gets 0$ \Comment{per query row}
\For{each key block $K_j, V_j$}
  \State $S_j \gets QK_j^{\top}/\sqrt{D}$; apply causal mask if enabled
  \State $m' \gets \max(m, \max(S_j))$
  \State $p \gets \exp(S_j - m')$
  \State $a \gets a \cdot \exp(m - m') + pV_j$; $\ell \gets \ell \cdot \exp(m - m') + \sum p$
  \State $m \gets m'$
\EndFor
\State \Return $a / \ell$
\end{algorithmic}
\end{algorithm}

\subsection{Greedy speculative decoding}
The draft model proposes $\gamma$ tokens with its own cache; the target verifies
$[\text{pending}, p_1, \dots, p_\gamma]$ in one forward pass, accepts the longest matching prefix and
appends its own token at the first disagreement. The emitted sequence is exactly the target's greedy
sequence. Cache bookkeeping is the delicate part: after verification the target drops the rejected
tail, and the draft drops one fewer position because it never consumed the last proposal; when every
proposal matched, the draft consumes that last token to stay level. Algorithm~\ref{alg:spec}
summarises the round. Tokenizers must agree on vocabulary and special tokens, which is asserted
before generation.

\begin{algorithm}[t]
\caption{One speculative round}
\label{alg:spec}
\begin{algorithmic}[1]
\Require pending token $t$, draft budget $\gamma$, caches $c_T, c_D$
\State $q \gets []$; $u \gets t$
\For{$i \gets 1$ to $\gamma$}
  \State $(c_D, \ell_D) \gets \textsc{Forward}_D([u], c_D)$; $u \gets \arg\max \ell_D$; append $u$ to $q$
\EndFor
\State $(c_T, L) \gets \textsc{Forward}_T([t, q_1, \dots, q_\gamma], c_T)$
\State $m \gets$ length of longest prefix with $q_i = \arg\max L_i$
\State \Return $q_{1..m} \parallel [\arg\max L_{m+1}]$, cropping $c_T$ by $\gamma-m$ and $c_D$ accordingly
\end{algorithmic}
\end{algorithm}

\section{Experimental Setup and Metrics}
All fine-tuning and evaluation runs use the T4 GPU available in Colab; local development and unit
tests run on CPU with tiny random Qwen2 models. Generation is greedy with at most 128 new tokens.
Classification is reported per criterion type with macro-averaged precision, recall and F1, plus a
confusion matrix; evidence quality is reported with micro and macro precision, recall and F1, where
an empty gold set counts as correct only for an empty prediction. Decoding benchmarks report
latency split into prefill and decode, peak allocated memory and forward-call counts; attention
benchmarks compare against a full-matrix fp32 reference and PyTorch SDPA; speculative decoding
reports acceptance rate and speedup against the cached target on eight validation prompts.

\section{Results and Discussion}

\subsection{Fine-tuning and classification}
Figure~\ref{fig:curves} shows the training and validation loss: the best validation loss is reached
at epoch 2 ($0.1170$ versus $0.1349$ after epoch 1 and $0.1227$ after epoch 3), and that checkpoint
is the adapter evaluated below. The run took 697\,s on the T4 (train loss $0.1022$). Table~\ref{tab:classification} compares the base and fine-tuned
models on validation, and Table~\ref{tab:test} reports the single test pass.

The base model does not follow the task at all: on validation it never emits an inclusion label and
predicts only \texttt{not applicable} (68 of 206 records) or \texttt{excluded} (138), giving $1.5\%$
accuracy. The fine-tuned adapter follows the JSON contract on every record (strict parse $100\%$,
invalid $0\%$) and reaches $70.4\%$ validation and $71.4\%$ test accuracy, but it collapses the rare
labels: it never predicts \texttt{not included}, \texttt{not applicable} or \texttt{excluded}, so
those per-class F1 scores are zero and macro-F1 stays at $0.34$/$0.37$ while accuracy looks much
better. The gap between accuracy and macro-F1 is the main result to keep in mind: the model is
usable as a majority-class screening aid, not as a balanced four-way classifier.

Table~\ref{tab:classification} also contains the hand-written KV-cached decoder. On all 206
validation records it produced \textbf{byte-identical raw output} to \texttt{model.generate}
(identical accuracy, confusion matrix and every metric), which is the real-model equivalence check
for the cache implementation; only latency changes.

\begin{figure}[t]
\centering
\pendingfigure{figures/train_curves.pdf}{Training curves: generated by Colab job 1.}
\caption{Training and validation loss for the LoRA fine-tune.}
\label{fig:curves}
\end{figure}

\begin{table}[t]
\centering
\caption{Validation classification results. Precision/recall/F1 are macro averages over the six
legal labels; invalid is the share of unparseable outputs.}
\label{tab:classification}
\begin{tabular}{llcccc}
\toprule
Model & Decode & Prec. & Rec. & F1 & Invalid \\
\midrule
Base & hf & 0.005 & 0.181 & 0.009 & 0.0\% \\
Fine-tuned & hf & 0.358 & 0.337 & 0.338 & 0.0\% \\
Fine-tuned & cached & 0.358 & 0.337 & 0.338 & 0.0\% \\
\bottomrule
\end{tabular}
\end{table}

\begin{table}[t]
\centering
\caption{Test-split results, evaluated once with \texttt{-{}-final}.}
\label{tab:test}
\begin{tabular}{lcccc}
\toprule
Model & Macro F1 & Evidence micro F1 & Invalid rate & $n$ \\
\midrule
Base & 0.053 & 0.050 & 0.0\% & 206 \\
Fine-tuned & 0.373 & 0.543 & 0.0\% & 206 \\
\bottomrule
\end{tabular}
\end{table}

Figure~\ref{fig:confusion} shows the overall confusion matrix (per-type matrices are in
\texttt{report/figures}). Inclusion and exclusion criteria are kept separate because their label sets
are disjoint. Errors are dominated by the majority class: 28 of the 66 gold
\texttt{not enough information} records on test are predicted \texttt{not excluded}, 12 of 15
\texttt{not applicable} records become \texttt{not excluded}, and 9 of 27 \texttt{included} records
become \texttt{not enough information}; the reverse direction appears too, with 7 gold
\texttt{not included} records predicted \texttt{not enough information}. No test record is predicted
\texttt{excluded}, so that column of the matrix is empty. The pattern is consistent with a model that
learned the majority label of each criterion type and the \texttt{not enough information} fallback,
rather than the boundary between ``satisfied'', ``not satisfied'' and ``inapplicable''.

\begin{figure}[t]
\centering
\pendingfigure{figures/val_finetuned_hf_confusion_overall.pdf}{Confusion matrices: generated by Colab job 3.}
\caption{Confusion matrices for the fine-tuned model on validation data.}
\label{fig:confusion}
\end{figure}

\subsection{Evidence quality and error analysis}
Table~\ref{tab:evidence} reports evidence precision, recall and F1 under the empty-set rule; 116 of
the 206 validation records have empty gold evidence, which is why macro-F1 ($0.76$) is much higher
than micro-F1 ($0.57$): the model often cites something when the expert cited nothing. The
fine-tuned model is nevertheless far better than the base model, which emits unusable citations
(micro-F1 $0.10$, and 69 empty predictions against 116 empty gold sets). On test the same pattern
holds (micro-F1 $0.54$, macro-F1 $0.80$). The error-analysis file \texttt{results/error\_cases.md}
contains 92 cases with the note, criterion, gold label and prediction; Table~\ref{tab:errors} groups
them by cause. Wrong labels dominate (56), followed by missing evidence (16), unsupported citations
(9), over-citation (8) and the \texttt{not applicable}/\texttt{not enough information} confusion
(3).

\begin{table}[t]
\centering
\caption{Evidence-retrieval metrics on validation data (empty-set rule; ``empty preds'' counts
records where the model cited nothing, against 116 records with empty gold evidence).}
\label{tab:evidence}
\begin{tabular}{lccc}
\toprule
Model & Micro F1 & Macro F1 & Empty preds \\
\midrule
Base & 0.101 & 0.167 & 69 \\
Fine-tuned & 0.569 & 0.761 & 137 \\
\bottomrule
\end{tabular}
\end{table}

\begin{table}[t]
\centering
\caption{Error categories in the test error analysis.}
\label{tab:errors}
\begin{tabular}{lcc}
\toprule
Category & Count & Example note \\
\midrule
Wrong label & 56 & ann.\ 161 \\
Missing evidence & 16 & ann.\ 179 \\
Unsupported citation & 9 & ann.\ 166 \\
Over-cited evidence & 8 & ann.\ 306 \\
App.\ vs.\ insufficient info & 3 & ann.\ 422 \\
\bottomrule
\end{tabular}
\end{table}

\subsection{Decoding benchmarks}
Table~\ref{tab:decoding} and Figure~\ref{fig:decoding} compare uncached and cached greedy decoding
over the $3\times3$ grid of input and output lengths. All nine grid points are token-identical, which
is the expected result for greedy decoding and confirms the hand-written loop. The cache pays off as
the prompt grows: at 128 input tokens the speedup is only $1.04$--$1.18\times$ because prefill
dominates and each cached step still launches a forward pass, while at 1{,}024 input tokens and 128
new tokens it reaches $5.11\times$ ($28.86$\,s versus $5.64$\,s). Across the grid the speedup ranges
from $1.04\times$ to $5.11\times$, and peak allocated memory grows with the cache up to $9.5$\,GB at
the longest setting.

\begin{table}[t]
\centering
\caption{Decoding latency (seconds, mean of repeated runs; three representative grid points of
nine, all token-identical).}
\label{tab:decoding}
\begin{tabular}{rrccc}
\toprule
Input & Output & Uncached & Cached & Speedup \\
\midrule
128 & 32 & 1.431 & 1.314 & 1.09$\times$ \\
512 & 128 & 16.804 & 5.545 & 3.03$\times$ \\
1024 & 128 & 28.855 & 5.644 & 5.11$\times$ \\
\bottomrule
\end{tabular}
\end{table}

\begin{figure}[t]
\centering
\pendingfigure{figures/bench_decoding.pdf}{Decoding latency and throughput: generated by job 5.}
\caption{Uncached versus KV-cached decoding.}
\label{fig:decoding}
\end{figure}

\subsection{Attention benchmarks}
Table~\ref{tab:attention} and Figure~\ref{fig:attention} compare the reference, the tiled kernel and
the available SDPA backends. On the T4 (sm\_75) PyTorch exposes the math and memory-efficient SDPA
backends; the FlashAttention kernel is reported unsupported because it requires sm80 or newer. The
tiled kernel agrees with the reference within $2.1\times10^{-3}$ maximum absolute error in fp16
(tolerance $10^{-2}$), so the online softmax is numerically correct. It is nevertheless much slower
in wall-clock time --- $0.285$\,s versus $0.005$\,s at 2{,}048 tokens --- because each key block is a
separate Python-level operation with no fused SRAM kernel, exactly the overhead FlashAttention was
designed to remove. Its purpose here is memory: at 2{,}048 tokens it peaks at $26$\,MiB versus
$158$\,MiB for the full-matrix reference and $336$\,MiB for SDPA math, because it never materialises
the $T\times T$ score matrix. The memory-efficient SDPA backend is both faster ($0.9$\,ms) and
smaller ($18$\,MiB), which is the honest comparison point.

\begin{table}[t]
\centering
\caption{Attention latency (seconds) and tiled-kernel maximum absolute error versus the fp32
reference (fp16 inputs; FlashAttention unsupported on sm\_75).}
\label{tab:attention}
\begin{tabular}{rcccc}
\toprule
Length & Reference & Tiled & SDPA (best) & Max error \\
\midrule
256 & 0.0006 & 0.0075 & 0.0002 & $2.1\times10^{-3}$ \\
512 & 0.0005 & 0.0237 & 0.0002 & $1.9\times10^{-3}$ \\
1024 & 0.0012 & 0.0903 & 0.0004 & $1.4\times10^{-3}$ \\
2048 & 0.0046 & 0.2848 & 0.0009 & $2.1\times10^{-3}$ \\
\bottomrule
\end{tabular}
\end{table}

\begin{figure}[t]
\centering
\pendingfigure{figures/bench_attention.pdf}{Attention latency and memory: generated by job 5.}
\caption{Attention variants against sequence length.}
\label{fig:attention}
\end{figure}

\subsection{Speculative decoding}
Table~\ref{tab:speculative} and Figure~\ref{fig:speculative} report acceptance rate, forward-call
counts and the change in latency against the cached target. The hand-written draft/verify loop is
lossless on every prompt (0 mismatches over eight prompts and 1{,}024 generated tokens), and the
independent Hugging Face \texttt{assistant\_model} check matches on the full 16-token prefix it
produces. Acceptance is high ($82\%$ at $\gamma=2$), so the draft and target mostly agree.

The speedup, however, is negative: all three settings are $1.5$--$1.6\times$ \emph{slower} than plain
cached decoding ($0.64$--$0.68\times$). The reason is hardware-specific: the 0.5B draft is only about
three times cheaper per forward than the 1.5B target, and our loop runs $\gamma$ sequential draft
steps per verification round, then only one target pass. On a T4 the draft steps plus the Python
loop overhead cost more than the saved target passes, even though acceptance is high. Batched or
fused drafting, a smaller draft, or a target/draft pair with a larger cost gap would be needed to
turn the acceptance rate into wall-clock speedup; we report the slowdown rather than hiding it.
$\gamma=2$ is the least bad setting, which is consistent with the draft cost dominating.

\begin{table}[t]
\centering
\caption{Speculative decoding on eight validation prompts (fixed 128-token budget; cached target
baseline 4.812\,s, $26.6$ tokens/s; every row token-identical).}
\label{tab:speculative}
\begin{tabular}{cccccc}
\toprule
$\gamma$ & Acceptance & Target calls & Draft calls & Latency (s) & Speedup \\
\midrule
2 & 0.820 & 395 & 1071 & 7.036 & 0.68$\times$ \\
4 & 0.698 & 280 & 1240 & 7.340 & 0.66$\times$ \\
6 & 0.632 & 224 & 1390 & 7.577 & 0.64$\times$ \\
\bottomrule
\end{tabular}
\end{table}

\begin{figure}[t]
\centering
\pendingfigure{figures/bench_speculative.pdf}{Speculative latency and acceptance: job 5.}
\caption{Speculative decoding versus the cached target.}
\label{fig:speculative}
\end{figure}

\subsection{Application}
Figure~\ref{fig:app} shows the screening tab and the live benchmark tab. \pending{insert the two
screenshots taken by the user and describe the first-load time on CPU and GPU.}

\begin{figure}[t]
\centering
\pendingfigure{figures/app_screening.png}{Screening tab screenshot (user).}
\pendingfigure{figures/app_benchmark.png}{Benchmark tab screenshot (user).}
\caption{The Gradio application.}
\label{fig:app}
\end{figure}

\section{Limitations}
The label distribution is highly imbalanced: only ten \texttt{excluded} and twenty-three
\texttt{not included} examples exist, so the corresponding per-class metrics are noisy and a single
test run cannot resolve small differences. The aggregation rule is a documented heuristic, not a
validated clinical decision procedure, and the model has no access to structured laboratory data
that is absent from the note. Evidence is extracted as sentence identifiers, so a decision can be
supported by a sentence that is topically related but not sufficient. Finally, all experiments use a
single seed and one hardware target; the inference benchmarks measure one GPU and should not be
generalised to other accelerators without repetition.

\section{Conclusion}
We fine-tuned a 1.5B instruction model for criterion-level trial screening with an auditable JSON
contract, implemented KV-cached decoding, tiled attention and speculative decoding from scratch, and
verified each against a reference before benchmarking it. The resulting application exposes
screening, a deterministic report with a persistent disclaimer, and a live comparison of the three
decoders. The fine-tune lifts validation accuracy from $1.5\%$ to $70.4\%$ and test accuracy from
$1.9\%$ to $71.4\%$, with macro-F1 limited to $0.37$ by the collapsed rare labels. On the inference
side, the hand-written KV cache is token-identical to the framework and up to $5.1\times$ faster,
the tiled kernel is numerically exact but slower than fused SDPA while using six times less memory,
and speculative decoding is lossless but $1.5\times$ slower on this GPU. Those three outcomes are
reported as measured: the exercise was to implement and verify the mechanisms, not to claim a
speedup that the hardware does not deliver.

\section*{AI-Use Appendix}
This project was built with heavy AI assistance, documented honestly in \texttt{docs/AI\_USE.md}.
The coding agent (OpenCode with a DeepSeek model) wrote the implementation and tests, ran the local
quality gate and smoke checks, and drafted this skeleton; the human author specified the
requirements, ran the Colab jobs, verified the results and owns every claim. No result in this paper
may be reported until it exists in \texttt{results/}.

\bibliographystyle{IEEEtran}
\bibliography{refs}

\end{document}
